\chapter{绪论}
%
\section{研究背景和意义}
\subsection{研究背景和意义}
% 大背景
随着新一轮科技革命与产业变革深入推进，以智能制造为核心的产业升级正在加速落地。在中国制造业从“规模扩张”向“质量提升与柔性生产”转型的背景下，机器人在制造与服务领域的应用正从“单一工艺、固定工位”的自动化，向“多品种、小批量、快速换型”的柔性智能化演进，工业机器人在3C电子、汽车零部件、金属加工等领域的应用需求持续增长。同时，服务机器人在物流、医疗康复、公共服务等场景也对“快速部署、快速适配、持续迭代”的能力提出了更高要求。由此带来的关键问题是：机器人系统不仅要“能做”，更要“能快速学会并稳定做好”，即在有限的时间与成本约束下完成新任务导入与可靠运行。以产业升级为牵引，企业对机器人提出了更强的工程诉求：不仅要实现高重复精度的自动运行，还要能够在短周期内完成新任务导入，并在工件几何误差、装夹偏差、摩擦与刚度变化等不确定因素存在时，保持稳定的作业质量与安全性。对于打磨抛光、装配插接、擦拭涂胶等典型任务，机械臂末端运动轨迹与接触交互过程高度耦合，轨迹获取与执行控制成为影响部署效率与任务成功率的关键环节。
% 获取轨迹的方法（补）--pid--pid改进--柔顺控制--神经网络（引用一些文献加上对应的图）
为提高轨迹跟踪精度与执行稳定性，工业机器人控制层面长期依赖较为成熟的经典控制方法。例如，基于误差反馈的 PID 控制因结构简单、易于整定、工程实践丰富，被广泛用于关节伺服和轨迹跟踪；在一定程度上，PID 能满足重复性任务的精度要求。但在面对强非线性、强耦合和接触扰动时，仅靠线性误差反馈往往难以兼顾精度与鲁棒性。与此同时，基于模型的计算力矩法、增益调度、前馈补偿等方法也被用于提升动态响应与跟踪性能。在涉及接触交互的任务中，“位置轨迹跟踪”本身并不足以保证安全与质量，控制目标需要从单纯的位置精度扩展到位置—力的协同控制。因此，阻抗控制/导纳控制等柔顺控制框架成为接触任务的重要基础：通过引入等效质量—弹簧—阻尼关系，使机械臂在接触过程中具有期望的柔顺动态响应，从而缓解力冲击并提升交互稳定性。近年来，随着传感器与算力发展，以及数据驱动方法的兴起，机器人轨迹示教与控制逐步从“纯模型/纯控制”向“控制与学习融合”演进。一类典型思路是利用神经网络学习系统建模动态、摩擦与接触效应，或学习从感知到控制的映射，从而在一定程度上降低精确建模与人工调参负担。但纯监督学习方法往往面临分布偏移与泛化不足：当执行状态偏离训练数据分布时，误差容易累积放大。相比之下，强化学习通过与环境交互优化策略，能够在不确定性与复杂约束条件下学习自适应行为，近年来在机器人控制、接触操作与技能学习方面显示出潜力。特别是将示教先验与强化学习结合（例如以示教轨迹为名义参考、由强化学习学习残差补偿或参数自适应），有望在保持轨迹结构与可解释性的同时，获得对环境变化的在线修正能力；再结合导纳/阻抗等柔顺控制作为执行层“安全骨架”，可以在接触任务中兼顾稳定性与学习效率。
% 指出问题
在上述技术演进背景下可以看到，现阶段关于机械臂的末端轨迹示教仍存在若干共性瓶颈，首先，从示教轨迹本身来看，示教数据往往以离散采样点形式记录，受示教者操作抖动、传感噪声与采样率限制影响，容易出现轨迹不平滑、速度剖面不合理、轨迹生成效率较低等问题。在不确定环境与接触任务中，“记录—回放”策略的局限性更为突出。示教得到的轨迹点序列通常隐含了特定工件几何、装夹位置与接触条件下的运动结果，当执行场景发生变化，末端轨迹即使在自由空间仍可能产生偏离，出现超限风险，造成设备磨损、工件损伤甚至安全隐患。在机器人与工件接触的阶段，系统会受到环境刚度突变、摩擦非线性、传感噪声与控制延迟等因素影响，易诱发振荡或颤振；同时，柔顺参数通常依赖人工整定，跨工况稳定性难以保证。因此，轨迹结构化表示与高效生成、在线自适应修正能力成为工业机器人轨迹示教研究的热点方向。
% 本文研究意义
本文围绕“示教轨迹如何获得—如何参数化—如何自适应优化—如何在接触中稳定执行”的主线，面向机械臂末端轨迹快速部署与接触任务稳定执行的工程需求，开展了基于强化学习的轨迹示教方法研究。首先，在示教获取层面，本文融合手眼标定与视觉伺服技术，实现机器人坐标系、视觉坐标系与工件坐标系的统一，并在视觉反馈下完成关键点定位与轨迹参考生成，从而提高示教效率并降低对人工经验的依赖。其次，在轨迹表示与优化层面，本文采用动态运动基元（DMP）对示教轨迹进行结构化参数化，实现轨迹的平滑重构与可调节的时空表达；进一步引入强化学习在执行过程中学习轨迹的自适应修正/补偿策略，在保持示教轨迹先验结构的基础上对环境不确定性与接触扰动进行在线调整，以缓解接触建立阶段的力突变与超限风险。最后，在执行控制层面，本文面向接触任务的力—位姿耦合特性，构建强化学习驱动的稳定接触控制策略，在柔顺控制框架下对接触过程进行自适应调节与稳定性约束设计，从而提升机械臂在接触过渡阶段的稳定性与作业安全性。通过上述研究，本文形成了一套贯穿示教采集、轨迹参数化、学习型自适应优化与接触稳定执行的统一方法框架，为复杂工况下机械臂末端轨迹示教与接触作业的快速部署与鲁棒执行提供支撑。
%
关于\LaTeX{}以及基于\LaTeX{}写作的好处不再赘述。\LaTeX{}的入门资料推荐文献\parencite{_g}以及文献\parencite{_c}。

这里主要是想推荐一种“学术生态”，即利用各种工具展开科研工作，以达到事半功倍的效果。需要用到以下软件：
\begin{enumerate}[topsep = 0 pt, itemsep= 0 pt, parsep=0pt, partopsep=0pt, leftmargin=44pt, itemindent=0pt, labelsep=6pt, label=(\arabic*)]
	\item 	参考文献管理软件zotero\cite{_m}。很多人使用过endnote，但其实zotero也非常强大，强烈推荐。可到b站观看Struggle with Me出品的视频教程\cite{_k}入门（或其他最新教程，刚开始不推荐使用插件，会增加学习难度）。zotero自带pdf阅读器，也可以设置为使用其他阅读器。在zotero可以打开文件所在位置，故不推荐更改zotero的文件系统（尤其不推荐使用zotfile插件，事实上各种五花八门的插件增加了复杂性，实际上没有带来太多便利性）。理论上只需要包含文献元数据信息的bib文件（可以手动一篇一篇文章地收集）即可使用此模板，因此模板不依赖于任何参考文献管理软件，endnote用户或不使用参考文献管理软件的用户可以忽略本文zotero部分的讲解。
	\item	可截图获取文献中公式的软件mathpix\cite{_h}。在阅读别人的论文时，很可能需要把文章中的公式抄下来放到自己的笔记中，方便以后组会报告甚至论文中使用，这时使用mathpix可直接截图获取\LaTeX{}源码，非常方便。随着mathpix的使用成本越来越高，免费次数越来越少，2023起已经不再推荐。目前开源/免费的替代工具为：\href{https://www.simpletex.cn/}{SimpleTex}和\href{https://p2t.breezedeus.com/}{Pix2Tex}。现在（2025年）SimpleTex网页不收费但客户端早已经收费，Pix2Tex虽然一般般但是好过没有，还是不错的选择，毕竟一般来说公式不会太复杂。Pix2Tex也一直在改进，可以本地部署也可以在线使用，MacOS还有桌面应用程序。
	\item	TeXlive202x、TeXstudio（2022起推荐vscode），相当于开发环境和IDE。本模板是基于TeX的发行版TeXlive202x和编辑器TeXstudio进行的，百度这两个关键字分别安装。关于TeXstudio的使用（快捷键等）可另行查找资料。模板还支持更多ide，更多编译方式见GitHub首页readme.md。若在其他窗口打开了编译生成的pdf文件，记得关掉再编译，否则报错。TeXstudio的设置见第二章。vscode的设置见github讨论区的\href{https://github.com/mengchaoheng/SCUT_thesis/discussions/6}{vscode配置}。
\end{enumerate}

本文的章节安排如下：

第一章，绪论。

第二章，模板简介。主要介绍各文件的内容。

第三章，常用环境。介绍论文写作中常用的环境，包括：图、表、公式、定理。基本涵盖了常用的命令。

%第三章，参考文献设置。本模板对旧版的改动主要是参考文献部分，本章将简单参考文献设置以及
%编译选项的设置等等。


